\par\noindent\textbf{(1)} 线性温度梯度下 \(\delta n_{\boldsymbol k}=-\tau v_x(\partial n_B/\partial T)\partial_xT\)。热流 \(j_{Qx}=\int d^3k/(2\pi)^3\,\hbar\omega\,v_x\delta n\)，球对称角平均 \(\langle v_x^2\rangle=v^2/3\)，故 \(\kappa=(1/3)\int d\omega\,g(\omega)\hbar\omega\,v^2\tau\,\partial_Tn_B\)。
\par\noindent\textbf{(2)} 若边界散射主导且 \(\tau\) 近似常数、\(T\ll\Theta_D\)，Debye热容 \(C_V\propto T^3\)，动理学公式简化为 \(\kappa\simeq C_Vv^2\tau/3\propto T^3\)。若平均自由程 \(\ell=v\tau\)，也写 \(\kappa\simeq C_Vv\ell/3\)。
\par\noindent\textbf{(3)} 普通声子碰撞在第一Brillouin区内守恒晶格准动量，而Umklapp过程满足 \(\boldsymbol k_1+\boldsymbol k_2=\boldsymbol k_3+\boldsymbol G\)，\(\boldsymbol G\ne0\) 为倒格矢，可将总声子准动量转移给晶格，从而松弛热流。单纯能量守恒不保证热流不衰减。
\par\noindent\textbf{(4)} 在 \(T\gg\Theta_D\) 时比热趋Dulong–Petit常数，若Umklapp率近似 \(\tau_U^{-1}\propto T\)，则 \(\kappa\propto T^{-1}\)。这依赖寿命模型；强无序、局域振动、二维水动力声子输运或接近熔融时不能无条件使用此反比律。
